\section{Order on \texorpdfstring{$\mathbb{N}$}{N}}

\begin{topicbox}{Strict and Non-Strict Order}
\begin{exposition}
Order on a Peano system is extracted from addition. Since the chapter is using
the distinguished first element $1$ rather than an additive zero, the strict
order is the cleaner primitive notion: $a<b$ means that $b$ is reached from
$a$ by adding a nontrivial amount. The non-strict order is then obtained by
allowing equality in addition to strict increase.
\end{exposition}

\begin{definitionbox}{Definition (Strict Less Than on a Peano System)}
\begin{definition}[Strict Less Than on a Peano System]
\label{def:strict-less-than-on-peano-system}
Let $(P,S,1)$ be a Peano system, and let $a,b \in P$. One writes
\[
a<b
\]
exactly when there exists $c \in P$ such that
\[
b=a+c.
\]
\end{definition}
\end{definitionbox}

\begin{remark*}[Standard quantified statement]
\[
\begin{aligned}
&\text{Let } (P,S,1) \text{ be a Peano system and let } a,b \in P.\\
&a<b
\Longleftrightarrow
\exists c \in P\;(b=a+c).
\end{aligned}
\]
\end{remark*}

\begin{remark*}[Negated quantified statement]
\[
\begin{aligned}
&\text{Let } (P,S,1) \text{ be a Peano system and let } a,b \in P.\\
&a \not< b
\Longleftrightarrow
\forall c \in P\;(b \neq a+c).
\end{aligned}
\]
\end{remark*}

\begin{remark*}[Failure modes]
Strict inequality fails exactly when no additive witness in $P$ carries $a$ to
$b$. In that case, every nontrivial additive step from $a$ misses $b$.
\end{remark*}

\begin{remark*}[Failure mode decomposition]
\[
\begin{aligned}
&\text{Let } (P,S,1) \text{ be a Peano system and let } a,b \in P.\\
&a \not< b
\Longleftrightarrow
\underbrace{\forall c \in P\;(b \neq a+c)}_{\text{every proposed additive witness fails}}.
\end{aligned}
\]
\end{remark*}

\begin{remark*}[Interpretation]
The strict order records whether $b$ lies genuinely ahead of $a$ in the
counting structure. The witness $c$ is the amount of additive displacement
needed to move from $a$ to $b$.
\end{remark*}

\begin{dependencies}
\begin{itemize}
  \item \hyperref[def:addition-on-peano-system]{Addition on a Peano System}
\end{itemize}
\end{dependencies}

\begin{definitionbox}{Definition (Less Than or Equal on a Peano System)}
\begin{definition}[Less Than or Equal on a Peano System]
\label{def:less-than-or-equal-on-peano-system}
Let $(P,S,1)$ be a Peano system, and let $a,b \in P$. One writes
\[
a \le b
\]
exactly when
\[
a<b
\qquad\text{or}\qquad
a=b.
\]
\end{definition}
\end{definitionbox}

\begin{remark*}[Standard quantified statement]
\[
\begin{aligned}
&\text{Let } (P,S,1) \text{ be a Peano system and let } a,b \in P.\\
&a \le b
\Longleftrightarrow
\bigl(a<b \lor a=b\bigr).
\end{aligned}
\]
\end{remark*}

\begin{remark*}[Negated quantified statement]
\[
\begin{aligned}
&\text{Let } (P,S,1) \text{ be a Peano system and let } a,b \in P.\\
&a \not\le b
\Longleftrightarrow
\bigl(a \not< b \land a \neq b\bigr).
\end{aligned}
\]
\end{remark*}

\begin{remark*}[Failure modes]
The non-strict order fails in exactly two simultaneous ways: $a$ is not
strictly below $b$, and $a$ is not equal to $b$. Thus $b$ neither lies ahead
of $a$ nor coincides with it.
\end{remark*}

\begin{remark*}[Failure mode decomposition]
\[
\begin{aligned}
&\text{Let } (P,S,1) \text{ be a Peano system and let } a,b \in P.\\
&a \not\le b
\Longleftrightarrow
\underbrace{a \not< b}_{\text{no strict additive witness}}
\land
\underbrace{a \neq b}_{\text{equality also fails}}.
\end{aligned}
\]
\end{remark*}

\begin{remark*}[Interpretation]
The relation $a \le b$ allows either equality or a genuine forward step from
$a$ to $b$. It is the additive order relation derived from the strict witness
definition.
\end{remark*}

\begin{dependencies}
\begin{itemize}
  \item \hyperref[def:strict-less-than-on-peano-system]{Strict Less Than on a Peano System}
\end{itemize}
\end{dependencies}
\end{topicbox}

\begin{topicbox}{Basic Order Laws}
\begin{exposition}
Once the order relations have been defined, the first task is to show that
they satisfy the expected structural laws: reflexivity, transitivity, and
antisymmetry. These are the properties that make the additive comparison
relation behave as an honest order rather than a merely suggestive notation.
\end{exposition}

\begin{theorembox}{Theorem (Order Is Reflexive on a Peano System)}
\begin{theorem}[Order Is Reflexive on a Peano System]
\label{thm:order-reflexive-on-peano-system}
Let $(P,S,1)$ be a Peano system. For every $a \in P$,
\[
a \le a.
\]
\hyperref[prf:order-reflexive-on-peano-system]{Go to proof.}
\end{theorem}
\end{theorembox}

\begin{remark*}[Standard quantified statement]
\[
\begin{aligned}
&\text{Let } (P,S,1) \text{ be a Peano system.}\\
&\forall a \in P\;(a \le a).
\end{aligned}
\]
\end{remark*}

\begin{remark*}[Interpretation]
Every element is comparable to itself. Reflexivity supplies the equality case
inside the non-strict order.
\end{remark*}

\begin{dependencies}
\begin{itemize}
  \item \hyperref[def:less-than-or-equal-on-peano-system]{Less Than or Equal on a Peano System}
\end{itemize}
\end{dependencies}

\begin{theorembox}{Theorem (Order Is Transitive on a Peano System)}
\begin{theorem}[Order Is Transitive on a Peano System]
\label{thm:order-transitive-on-peano-system}
Let $(P,S,1)$ be a Peano system. For every $a,b,c \in P$,
\[
\bigl(a \le b \land b \le c\bigr)
\Longrightarrow
a \le c.
\]
\hyperref[prf:order-transitive-on-peano-system]{Go to proof.}
\end{theorem}
\end{theorembox}

\begin{remark*}[Standard quantified statement]
\[
\begin{aligned}
&\text{Let } (P,S,1) \text{ be a Peano system.}\\
&\forall a,b,c \in P\;
\bigl((a \le b \land b \le c) \Longrightarrow a \le c\bigr).
\end{aligned}
\]
\end{remark*}

\begin{remark*}[Interpretation]
If $b$ lies at or beyond $a$, and $c$ lies at or beyond $b$, then $c$ also
lies at or beyond $a$. Transitivity expresses the consistency of additive
comparison along a chain.
\end{remark*}

\begin{dependencies}
\begin{itemize}
  \item \hyperref[def:strict-less-than-on-peano-system]{Strict Less Than on a Peano System}
  \item \hyperref[def:less-than-or-equal-on-peano-system]{Less Than or Equal on a Peano System}
  \item \hyperref[thm:addition-associative]{Addition Is Associative}
\end{itemize}
\end{dependencies}

\begin{theorembox}{Theorem (Order Is Antisymmetric on a Peano System)}
\begin{theorem}[Order Is Antisymmetric on a Peano System]
\label{thm:order-antisymmetric-on-peano-system}
Let $(P,S,1)$ be a Peano system. For every $a,b \in P$,
\[
\bigl(a \le b \land b \le a\bigr)
\Longrightarrow
a=b.
\]
\hyperref[prf:order-antisymmetric-on-peano-system]{Go to proof.}
\end{theorem}
\end{theorembox}

\begin{remark*}[Standard quantified statement]
\[
\begin{aligned}
&\text{Let } (P,S,1) \text{ be a Peano system.}\\
&\forall a,b \in P\;
\bigl((a \le b \land b \le a) \Longrightarrow a=b\bigr).
\end{aligned}
\]
\end{remark*}

\begin{remark*}[Contrapositive quantified statement]
\[
\begin{aligned}
&\text{Let } (P,S,1) \text{ be a Peano system.}\\
&\forall a,b \in P\;
\bigl(a \neq b \Longrightarrow \neg(a \le b \land b \le a)\bigr).
\end{aligned}
\]
\end{remark*}

\begin{remark*}[Interpretation]
Two distinct elements cannot each lie at or beyond the other. Antisymmetry
separates genuine equality from bidirectional comparability.
\end{remark*}

\begin{dependencies}
\begin{itemize}
  \item \hyperref[def:strict-less-than-on-peano-system]{Strict Less Than on a Peano System}
  \item \hyperref[def:less-than-or-equal-on-peano-system]{Less Than or Equal on a Peano System}
  \item \hyperref[thm:addition-left-cancellation]{Left Cancellation for Addition}
  \item \hyperref[thm:addition-associative]{Addition Is Associative}
  \item \hyperref[thm:no-object-is-its-own-successor]{No Object Is Its Own Successor}
\end{itemize}
\end{dependencies}
\end{topicbox}

\begin{topicbox}{Compatibility with Addition}
\begin{exposition}
The order on a Peano system is built from addition, so it must interact
coherently with addition. The next theorems show that adding the same element
to both sides preserves the comparison relation, regardless of whether the
common summand is placed on the right or on the left.
\end{exposition}

\begin{theorembox}{Theorem (Addition Preserves Order on the Right)}
\begin{theorem}[Addition Preserves Order on the Right]
\label{thm:addition-preserves-order-on-right}
Let $(P,S,1)$ be a Peano system. For every $a,b,c \in P$,
\[
a \le b
\Longrightarrow
a+c \le b+c.
\]
\hyperref[prf:addition-preserves-order-on-right]{Go to proof.}
\end{theorem}
\end{theorembox}

\begin{remark*}[Standard quantified statement]
\[
\begin{aligned}
&\text{Let } (P,S,1) \text{ be a Peano system.}\\
&\forall a,b,c \in P\;
\bigl(a \le b \Longrightarrow a+c \le b+c\bigr).
\end{aligned}
\]
\end{remark*}

\begin{remark*}[Interpretation]
Adding the same quantity to two comparable elements does not disturb their
order. The additive witness for $a \le b$ survives after both sides are shifted
by the same summand.
\end{remark*}

\begin{dependencies}
\begin{itemize}
  \item \hyperref[def:strict-less-than-on-peano-system]{Strict Less Than on a Peano System}
  \item \hyperref[def:less-than-or-equal-on-peano-system]{Less Than or Equal on a Peano System}
  \item \hyperref[thm:addition-associative]{Addition Is Associative}
\end{itemize}
\end{dependencies}

\begin{theorembox}{Theorem (Addition Preserves Order on the Left)}
\begin{theorem}[Addition Preserves Order on the Left]
\label{thm:addition-preserves-order-on-left}
Let $(P,S,1)$ be a Peano system. For every $a,b,c \in P$,
\[
a \le b
\Longrightarrow
c+a \le c+b.
\]
\hyperref[prf:addition-preserves-order-on-left]{Go to proof.}
\end{theorem}
\end{theorembox}

\begin{remark*}[Standard quantified statement]
\[
\begin{aligned}
&\text{Let } (P,S,1) \text{ be a Peano system.}\\
&\forall a,b,c \in P\;
\bigl(a \le b \Longrightarrow c+a \le c+b\bigr).
\end{aligned}
\]
\end{remark*}

\begin{remark*}[Interpretation]
The same order preservation holds when the common summand is added on the
left. This is the left-handed form of the additive compatibility law.
\end{remark*}

\begin{dependencies}
\begin{itemize}
  \item \hyperref[thm:addition-preserves-order-on-right]{Addition Preserves Order on the Right}
  \item \hyperref[thm:addition-commutative]{Addition Is Commutative}
\end{itemize}
\end{dependencies}
\end{topicbox}

\begin{topicbox}{Compatibility with Multiplication}
\begin{exposition}
Multiplication by a natural number is repeated addition by a positive amount.
It therefore preserves strict comparisons, and because every multiplier lies
at or beyond $1$, it also reflects strict comparisons.
\end{exposition}

\begin{theorembox}{Theorem (Multiplication Preserves and Reflects Strict Order)}
\begin{theorem}[Multiplication Preserves and Reflects Strict Order]
\label{thm:multiplication-preserves-and-reflects-strict-order}
Let $(P,S,1)$ be a Peano system. For every $k,m,n \in P$,
\[
m<n
\Longleftrightarrow
k\cdot m<k\cdot n.
\]
\hyperref[prf:multiplication-preserves-and-reflects-strict-order]{Go to proof.}
\end{theorem}
\end{theorembox}

\begin{remark*}[Standard quantified statement]
\[
\begin{aligned}
&\text{Let } (P,S,1) \text{ be a Peano system.}\\
&\forall k,m,n \in P\;\bigl(m<n \Longleftrightarrow k\cdot m<k\cdot n\bigr).
\end{aligned}
\]
\end{remark*}

\begin{remark*}[Interpretation]
Multiplication by a fixed natural number is an order embedding of the natural
numbers into themselves. It shifts strict inequalities forward without
collapsing distinct values.
\end{remark*}

\begin{remark*}[Source alignment]
This is the house-style multiplication-order comparison theorem corresponding
to Feferman's order laws for products \citep{FefermanNumberSystems1964}.
\end{remark*}

\begin{dependencies}
\begin{itemize}
  \item \hyperref[def:strict-less-than-on-peano-system]{Strict Less Than on a Peano System}
  \item \hyperref[def:multiplication-on-peano-system]{Multiplication on a Peano System}
  \item \hyperref[thm:multiplication-distributes-over-addition]{Multiplication Distributes Over Addition}
  \item \hyperref[thm:multiplication-commutative]{Multiplication Is Commutative}
  \item \hyperref[thm:addition-preserves-order-on-right]{Addition Preserves Order on the Right}
  \item \hyperref[thm:addition-left-cancellation]{Left Cancellation for Addition}
\end{itemize}
\end{dependencies}
\end{topicbox}

\begin{topicbox}{Successor Characterization and Trichotomy}
\begin{exposition}
Strict inequality can be recast in terms of stepping one successor beyond the
smaller element. This successor reformulation sharpens the order relation and
leads to trichotomy, which states that any two elements are linearly
comparable.
\end{exposition}

\begin{theorembox}{Theorem (Strict Order Successor Characterization)}
\begin{theorem}[Strict Order Successor Characterization]
\label{thm:strict-order-successor-characterization}
Let $(P,S,1)$ be a Peano system. For every $a,b \in P$,
\[
a<b
\Longleftrightarrow
S(a) \le b.
\]
\hyperref[prf:strict-order-successor-characterization]{Go to proof.}
\end{theorem}
\end{theorembox}

\begin{remark*}[Standard quantified statement]
\[
\begin{aligned}
&\text{Let } (P,S,1) \text{ be a Peano system.}\\
&\forall a,b \in P\;
\bigl(a<b \Longleftrightarrow S(a) \le b\bigr).
\end{aligned}
\]
\end{remark*}

\begin{remark*}[Interpretation]
To say that $b$ is strictly larger than $a$ is to say that $b$ lies at or
beyond the first successor step past $a$. This theorem converts strict order
into a non-strict comparison involving successor.
\end{remark*}

\begin{dependencies}
\begin{itemize}
  \item \hyperref[def:strict-less-than-on-peano-system]{Strict Less Than on a Peano System}
  \item \hyperref[def:less-than-or-equal-on-peano-system]{Less Than or Equal on a Peano System}
  \item \hyperref[thm:addition-with-one]{Addition With $1$}
  \item \hyperref[thm:addition-successor-on-right]{Addition Successor on the Right}
  \item \hyperref[thm:addition-associative]{Addition Is Associative}
  \item \hyperref[thm:addition-commutative]{Addition Is Commutative}
  \item \hyperref[thm:every-element-is-one-or-a-successor]{Every Element Is Either $1$ or a Successor}
\end{itemize}
\end{dependencies}

\begin{theorembox}{Theorem (Trichotomy on a Peano System)}
\begin{theorem}[Trichotomy on a Peano System]
\label{thm:trichotomy-on-peano-system}
Let $(P,S,1)$ be a Peano system. For every $a,b \in P$, exactly one of the
following holds:
\[
a<b,\qquad a=b,\qquad b<a.
\]
\hyperref[prf:trichotomy-on-peano-system]{Go to proof.}
\end{theorem}
\end{theorembox}

\begin{remark*}[Standard quantified statement]
\[
\begin{aligned}
&\text{Let } (P,S,1) \text{ be a Peano system.}\\
&\forall a,b \in P\;
\Bigl(
(a<b \lor a=b \lor b<a) \land {}\\
&\qquad\neg(a<b \land a=b) \land
\neg(a<b \land b<a) \land
\neg(a=b \land b<a)
\Bigr).
\end{aligned}
\]
\end{remark*}

\begin{remark*}[Interpretation]
Any two elements of a Peano system are linearly comparable: one lies below the
other, or they coincide. Trichotomy is the theorem that upgrades the additive
comparison relation from a partial-order structure to a total order.
\end{remark*}

\begin{dependencies}
\begin{itemize}
  \item \hyperref[def:strict-less-than-on-peano-system]{Strict Less Than on a Peano System}
  \item \hyperref[def:less-than-or-equal-on-peano-system]{Less Than or Equal on a Peano System}
  \item \hyperref[thm:order-reflexive-on-peano-system]{Order Is Reflexive on a Peano System}
  \item \hyperref[thm:order-transitive-on-peano-system]{Order Is Transitive on a Peano System}
  \item \hyperref[thm:order-antisymmetric-on-peano-system]{Order Is Antisymmetric on a Peano System}
  \item \hyperref[thm:strict-order-successor-characterization]{Strict Order Successor Characterization}
  \item \hyperref[thm:induction-principle-for-peano-system]{Induction Principle for a Peano System}
\end{itemize}
\end{dependencies}
\end{topicbox}
